\section{模型选择与评价}
\label{sec:ensemble-evaluation}

\subsection{方法比较}
\label{sec:ensemble-method-comparison}

表\ref{tab:ensemble-method-navigation}按成员来源、依赖和组合规则比较本章方法。相同的“100个成员”在不同方法中不是同一计算单位：Bagging的一棵深树负责完整预测，Boosting的一棵浅树只提供增量，Stacking的一个成员甚至可能是已经包含许多树的随机森林。容量与成本必须沿完整训练和预测路径衡量。

\begin{table*}[htbp]
  \centering
  \small
  \begin{tabularx}{\linewidth}{@{}>{\raggedright\arraybackslash}p{23mm}*{6}{>{\raggedright\arraybackslash}X}@{}}
    \toprule
    方法 & 成员来源 & 成员间依赖 & 组合规则 & 训练并行性 & 推理路径 & 主要失效来源 \\
    \midrule
    Bagging & 重采样或随机权重下的同类成员 & 给定原训练集后可独立生成 & 平均、投票或固定加权 & 成员间可并行 & 通常运行全部成员 & 成员稳定或共同偏差占主导 \\
    随机森林／Extra Trees & 样本、特征与阈值的随机树 & 树之间无序贯依赖 & 平均或投票 & 树间可并行 & 运行全部或选定树 & 随机化削弱单树，或相关性仍高 \\
    AdaBoost & 当前样本权重下的弱分类器 & 第$t$轮依赖前轮权重 & 训练中确定的加权投票 & 轮次串行 & 累加全部或截断成员 & 噪声记录持续获得大权重 \\
    梯度提升树 & 当前损失梯度下的回归树 & 第$t$轮依赖$F_{t-1}$ & 收缩后的加法得分 & 轮次串行，单轮可并行 & 顺序累加树输出 & 损失、轮数、树复杂度与采样失配 \\
    Stacking & 同质或异质候选模型 & 基成员可独立，元层依赖折外预测 & 学习得到的元模型 & 折与成员可部分并行 & 运行基成员后再运行元模型 & 元特征泄漏、成员冗余或元层过拟合 \\
    动态选择／紧凑化 & 已训练成员池或教师集成 & 依赖能力估计或教师输出 & 输入相关选择或新学生模型 & 依具体训练流程而定 & 选择部分成员，或只运行学生 & 路由失准、覆盖不足或逼近误差 \\
    \bottomrule
  \end{tabularx}
  \caption{集成方法的属性导航。表中成本描述方法结构，不构成固定速度排序；实际结果还受成员模型、实现、硬件、数据规模与调参预算影响。}
  \label{tab:ensemble-method-navigation}
\end{table*}

方法比较还要固定组合输出。回归的均值、分位数和分布预测不能共用一个误差口径；分类的硬标签、排序得分和概率也需要不同证据。一个集成若只改善准确率，不能据此声称概率校准或尾部代价同步改善。具体指标沿用表\ref{tab:regression-evaluation-protocol}和表\ref{tab:classification-evaluation-protocol}，本节补充集成特有的数据隔离与资源记录。

\subsection{数据隔离与泄漏控制}
\label{sec:ensemble-data-isolation}

集成流程会产生多种“没有直接用于某次拟合的预测”，但它们的职责不同：
\begin{itemize}
  \item \term{袋外预测}来自未抽中当前样本的自助成员，用于Bagging内部诊断、置换重要性或相关构造。
  \item \term{折外预测}来自未使用当前折拟合的模型，用于训练Stacking元学习器，或为其他二阶段模型提供样本外特征。
  \item \term{能力估计预测}来自动态选择的独立开发数据，用于判断某个成员在当前输入附近是否可靠。
  \item \term{验证预测}用于选择成员数、树复杂度、学习率、组合约束、阈值与早停轮数。
  \item \term{测试预测}只在所有训练、选择、校准和部署规则冻结后产生，用于评价最终流程。
\end{itemize}
图\ref{fig:ensemble-data-isolation}把五类预测放回各自的数据路径。前四类仍处在模型开发阶段，允许影响规定范围内的诊断、二阶段训练或选择；测试预测位于流程冻结之后，只负责评价最终链路。区分这些职责，比笼统地把它们都称为“样本外预测”更能暴露泄漏。

\begin{figure*}[htbp]
  \centering
  % Data roles in ensemble development and final evaluation.
% Schematic only; no dataset sizes or empirical quantities are encoded.
% Canvas: 164 mm x 82 mm. Dependencies: project palette and TikZ.
\begin{tikzpicture}[
  x=1mm,y=1mm,
  font=\footnotesize,
  text=BookInk,
  source/.style={
    rectangle,
    draw=FigureInput,
    fill=FigureInput!8,
    line width=0.7pt,
    minimum width=43mm,
    minimum height=9mm,
    text width=39mm,
    align=center,
    inner sep=1mm
  },
  condition/.style={
    rectangle,
    draw=FigureLoss,
    fill=FigureLoss!9,
    line width=0.7pt,
    minimum width=52mm,
    minimum height=9mm,
    text width=48mm,
    align=center,
    inner sep=1mm
  },
  purpose/.style={
    rectangle,
    draw=FigureOutput,
    fill=FigureOutput!6,
    line width=0.8pt,
    minimum width=45mm,
    minimum height=9mm,
    text width=41mm,
    align=center,
    inner sep=1mm
  },
  final/.style={
    source,
    line width=0.9pt
  },
  final purpose/.style={
    purpose,
    line width=0.9pt
  },
  flow/.style={
    -{Latex[length=2.2mm,width=1.5mm]},
    draw=BookMuted,
    line width=0.8pt
  }
]
  \path[use as bounding box] (0,0) rectangle (164,82);

  \node[font=\heiti\small,text=FigureInput] at (25,78) {预测或数据来源};
  \node[font=\heiti\small,text=FigureLoss] at (82,78) {保持隔离的条件};
  \node[font=\heiti\small,text=FigureOutput] at (139,78) {在开发流程中的用途};

  \node[source] (oob-source) at (25,65) {袋外预测};
  \node[condition] (oob-condition) at (82,65)
    {成员$m$未使用样本$i$训练};
  \node[purpose] (oob-purpose) at (139,65)
    {Bagging内部诊断与构造};

  \node[source] (oof-source) at (25,51) {折外预测};
  \node[condition] (oof-condition) at (82,51)
    {折内模型未使用当前留出折};
  \node[purpose] (oof-purpose) at (139,51)
    {训练Stacking元学习器};

  \node[source] (competence-source) at (25,37) {能力估计预测};
  \node[condition] (competence-condition) at (82,37)
    {能力数据不参与成员训练};
  \node[purpose] (competence-purpose) at (139,37)
    {估计局部能力并路由};

  \node[source] (validation-source) at (25,23) {验证预测};
  \node[condition] (validation-condition) at (82,23)
    {与相应候选的拟合路径隔离};
  \node[purpose] (validation-purpose) at (139,23)
    {调参、筛选、校准与早停};

  \draw[BookRule,line width=0.8pt] (1,14)--(163,14);
  \node[fill=white,inner sep=1mm,font=\heiti\small,text=BookInk] at (82,14)
    {冻结全部训练、选择、校准与部署规则};

  \node[final] (test-source) at (25,6) {测试预测};
  \node[condition,line width=0.9pt]
    (test-condition) at (82,6)
    {测试样本未进入任何开发决策};
  \node[final purpose] (test-purpose) at (139,6)
    {评价最终完整流程};

  \foreach \prefix in {oob,oof,competence,validation,test} {
    \draw[flow] (\prefix-source)--(\prefix-condition);
    \draw[flow] (\prefix-condition)--(\prefix-purpose);
  }
\end{tikzpicture}

  \caption{集成学习中的数据隔离与预测职责。每一行从预测或数据来源出发，经由相应隔离条件到达允许的用途；水平分隔线表示训练、调参、筛选、校准和部署规则已经冻结。测试数据位于分隔线之后，不能沿开发路径回流。}
  \label{fig:ensemble-data-isolation}
\end{figure*}

这些预测都可能在局部意义上“样本外”，却不能互相替代。折外矩阵参与元学习器训练，因而不是最终测试证据；袋外误差若反复用于选择随机森林参数，也已经参与模型开发；动态能力数据若同时训练成员，会高估局部能力。

泄漏还可能发生在模型拟合之前。缺失值填补、标准化、特征选择、目标编码和降维若利用了当前留出折，就会把信息带入折外预测。基成员调参、成员池筛选、元学习器调参和概率校准也要放在正确的数据层级。需要同时比较候选成员和组合规则时，可以在外层保留最终测试集，在内层训练折中完成全部预处理、调参及折外特征构造；时间、患者、用户或设备等依赖单位则按部署边界分组划分。

数据隔离不是要求为每一步机械切出一份新数据。样本有限时，可以通过嵌套交叉验证或交叉拟合复用开发样本，同时保持当前预测来自未用该样本拟合的路径。应记录每个预测由哪些样本训练的哪些模型产生，使元训练、选择与最终评价的角色可追溯。

\subsection{预测质量与资源成本}
\label{sec:ensemble-quality-resource}

集成结果至少要与三个基线比较：相同任务下的合理单模型、未简化的完整集成，以及满足同一资源限制的候选方法。只比较集成与其中最差成员，会夸大组合价值；只比较最终分数，又无法判断收益来自成员质量提升、错误依赖变化还是组合规则。

对Bagging，可以同时报告代表性单成员损失、成员预测相关或双错率、集成损失以及损失随成员数的变化。对Boosting，应区分训练损失、验证轨迹和最终测试结果，并记录早停轮数。对Stacking，应比较各成员的折外表现、折外预测依赖、元模型的新增验证收益和全链路测试结果。多样性指标是诊断量，不应脱离任务损失单独优化。

概率输出需要额外检查。平均若干过度自信的概率，不保证得到可靠概率；AdaBoost和梯度提升的累计得分也需要经过适当链接才能解释为概率。若使用后处理校准，校准器只能接收未用于拟合该预测的数据，并要把它作为部署链路的一部分，在最终测试集上报告对数损失、Brier分数、可靠性及任务决策损失。

资源比较应记录完整过程，而不是只测单个成员。训练侧至少包括总计算量、墙钟时间、峰值内存、并行资源和调参次数；部署侧包括序列化模型大小、加载内存、单样本延迟、批量吞吐及预处理和路由开销。Bagging的成员训练可以并行，但推理通常运行多个成员；Boosting轮次串行，树内统计可并行；Stacking既要运行基成员，又要运行元模型。静态选择、动态选择和紧凑化是否真正节省资源，只能由端到端测量回答。

\subsection{从问题条件到候选方法}
\label{sec:ensemble-method-selection}

当单模型对训练样本、随机种子或局部分裂明显敏感，并且可以承担多个成员的存储与预测成本时，Bagging是直接候选。树成员反复依赖少数强特征时，可以进一步比较随机森林；若阈值搜索成本较高且可以接受更强的随机化，再比较Extra Trees。是否采用这些方法，应由固定任务上的单成员变化、成员依赖和组合收益共同支持。

当任务损失能够提供有意义的逐轮修正，并且训练流程可以承担序贯依赖时，可以考虑Boosting。AdaBoost适合检查重加权分类的机制；一般损失与树成员则进入GBDT。XGBoost、LightGBM和CatBoost的选择不应只依据名称：需要分别判断是否重视正则化二阶树目标与稀疏路径、直方图及大规模稀疏统计、或类别特征的有序统计，并在当前硬件和搜索预算下实测。

当候选模型已经产生互补的样本外预测，而且可以承担交叉验证与多模型部署成本时，Stacking提供学习式组合。数据较少时，固定留出会压缩元训练样本，完整折外构造通常更充分；成员高度相关时，受约束线性组合比高容量元模型更容易控制。Super Learner的理论比较范围由候选库、损失和组合类共同决定，不能替代有限样本验证。

成员池超过部署预算时，再判断要保留什么对象。希望所有输入走同一条较短路径时采用静态成员选择；有可靠局部能力数据且路由开销可控时考虑动态选择；完整集成质量合格但任何成员子集都达不到预算时，可以训练紧凑模型。以上条件只负责缩小候选范围，最终选择仍来自相同数据制度、调参预算和资源口径下的新样本证据。
